\subsection{Procesamiento}
\label{sec:prueba_p}
\textbf{Pregunta y propósito.} \textquestiondown{}Conserva DME su valor al crecer los archivos con iguales obligaciones y entregas? Una tarea permite cotejar las entregas de entrada y final.

\begin{table}[htbp]\centering\small
\caption{Configuraciones del contraste de procesamiento}
\begin{tabularx}{\linewidth}{@{}P{1.4cm}YY@{}}
\toprule Caso & Configuración & Por qué se incluye\\\midrule
P01 & Una tarea en W1, una entrada externa de 10 MiB y un final de 10 MiB en el coordinador. & Base de dos obligaciones, sin intermedios.\\
P02 & La entrada y el final miden 20 MiB. La tarea, el worker y el plan se conservan. & Duplica ambos tamaños para comprobar la escala por bytes, no por número de archivos.\\
\bottomrule
\end{tabularx}
\end{table}

\textbf{Factor y controles.} Varía el tamaño uniforme de entrada y final (10 o 20 MiB); se mantienen una tarea por job, W1, destino, recursos y política de entrega. W2 sigue disponible. El nivel de 10 MiB coincide con la referencia del pipeline.

\textbf{Comprobación esperada.} Entrada y final llevan requerido y observado de 20 a 40 MiB; se esperan dos entregas y DME igual a uno en ambos casos. Se cotejan identidad y tamaño; el tiempo de DAGMan se registra por separado.

\textbf{Resultado y decisión.} P01 y P02 corroboraron dos entregas; requerido y observado pasaron de 20 a 40 MiB y DME se mantuvo en 1. La razón se conserva, aunque los volúmenes absolutos distinguen las cargas.

\subsection{Pipeline}
\label{sec:prueba_l}
\textbf{Pregunta y propósito.} Con flujo y asignación fijos, \textquestiondown{}qué agrupación entrega menos volumen? El pipeline compara una o dos dependencias internas; I40 aumenta un intermedio con el grafo fijo.

\begin{table}[htbp]\centering\small
\caption{Configuraciones del contraste de pipeline}
\begin{tabularx}{\linewidth}{@{}P{1.35cm}YY@{}}
\toprule Caso & Configuración & Por qué se incluye\\\midrule
L01 & $T_1$, $T_2$ y $T_3$ en tres jobs separados en W1. Entrada, dos intermedios y final de 10 MiB. & Base sin agrupación; expone ambas fronteras entre jobs.\\
GL2 & $T_1$ y $T_2$ en un job y $T_3$ en otro. Todas las tareas en W1. & Hace interna la primera dependencia y conserva la segunda frontera.\\
GL3 & Las tres tareas se ejecutan secuencialmente en un job en W1. & Elimina ambas fronteras y permite examinar los retornos declarados.\\
I40 & Mismo plan que L01, con $f_{12}$ de 40 MiB y los otros tres archivos de 10 MiB. & Aumenta cuatro veces un intermedio, con tareas y entregas previstas constantes.\\
\bottomrule
\end{tabularx}
\end{table}

\textbf{Factor y controles.} La agrupación cambia los jobs científicos (tres, dos, uno) con tamaños, tareas, dependencias, W1 y entrega constantes. L01 e I40 comparan tamaño en tres jobs.

\textbf{Comprobación esperada.} Los cuatro casos requieren 20 MiB por entrada y final. Con los retornos, L01, GL2 y GL3 implican seis, cinco y cuatro entregas de 10 MiB. En I40, retorno y reenvío añaden $2(40-10)=60$ MiB frente a L01; se cotejan con el plan y los registros.


\noindent\textbf{Resultado y decisión del pipeline.} L01, GL2 y GL3 observaron 60, 50 y 40 MiB; con 20 MiB requeridos, DME fue $1/3$, $2/5$ y $1/2$. Frente a L01, GL2 redujo el observado 16,7\,\% y GL3 33,3\,\%; GL3 tuvo el menor observado de las tres agrupaciones históricas con la misma política de entrega. Los dos retornos declarados explican el exceso.

En I40, el primer intermedio creció de 10 a 40 MiB. Retorno y reenvío añadieron 30 MiB cada uno: el observado pasó de 60 a 120 MiB. Productor y consumidor permanecen en W1; requerido siguió en 20 MiB y DME bajó de $1/3$ a $1/6$. El efecto depende del archivo y sus entregas; una DME menor no indica un fallo del cálculo.

\subsection{Agregación}
\label{sec:prueba_a}
\textbf{Pregunta y propósito.} Si el agregador cambia de máquina con los mismos productores y archivos, \textquestiondown{}cuánto cambia la referencia y cuánto cambia el observado? Distingue nuevas obligaciones de transferencias del mecanismo.

\begin{table}[htbp]\centering\small
\caption{Configuraciones del contraste de agregación}
\begin{tabularx}{\linewidth}{@{}P{1.4cm}YY@{}}
\toprule Caso & Configuración & Por qué se incluye\\\midrule
A01 & Cuatro productores y un agregador en W1, en cinco jobs. Cuatro entradas y cuatro intermedios de 10 MiB, y un final de 40 MiB. & Coloca productores y consumidor en el mismo sitio. Permite separar los retornos del mecanismo de las obligaciones entre máquinas.\\
A02 & Los cuatro productores se mantienen en W1. El agregador pasa a W2. Se conservan los cinco jobs y los archivos. & Cambia solo la ubicación del agregador y muestra el requerido por los nuevos destinos.\\
\bottomrule
\end{tabularx}
\end{table}

\textbf{Factor y controles.} Se varía el worker del agregador (W1 o W2); se mantienen la asignación de productores, los jobs, las identidades de archivos y la entrega final al coordinador. El final mide 40 MiB, igual a la suma de los cuatro intermedios.

\textbf{Comprobación esperada.} Entradas y final suman 80 MiB en A01. En A02, los cuatro intermedios de 10 MiB retornan al coordinador y se entregan a W2, por lo que se añaden $4(10+10)=80$ MiB y la cota es 160 MiB. El observado es 160 MiB en ambos casos.

\textbf{Resultado y decisión.} A01 tiene una cota de 80 MiB y un observado de 160 MiB: DME es $1/2$. En A02 ambos son 160 MiB y DME es uno; los retornos intermedios son relevo obligatorio. La cota cambia con la ubicación; ambas configuraciones empatan en volumen observado.

\subsection{Distribución}
\label{sec:prueba_d}
\textbf{Pregunta y propósito.} \textquestiondown{}Puede la agrupación conservar las entregas cuando los consumidores necesitan archivos distintos? Se contrasta con el pipeline, que tiene dependencias internas.

\begin{table}[htbp]\centering\small
\caption{Configuraciones del contraste de distribución}
\begin{tabularx}{\linewidth}{@{}P{1.4cm}YY@{}}
\toprule Caso & Configuración & Por qué se incluye\\\midrule
D02 & Un divisor en W1 y cuatro consumidores en W2, en jobs separados. Entrada de 10 MiB y cuatro particiones y finales de 2,5 MiB cada uno. & Es una base con cuatro consumidores de archivos distintos y con una ubicación que exige entregar cada partición a W2.\\
GD4 & El divisor permanece separado en W1 y los cuatro consumidores se ejecutan en un job en W2. & Agrupa todos los consumidores; comprueba si conserva las cuatro entregas.\\
\bottomrule
\end{tabularx}
\end{table}

\textbf{Factor y controles.} Se varía el número de jobs de consumo (cuatro o uno), con ubicaciones, grafo y política de entrega fijos. Cada consumidor lee una partición y devuelve un final del mismo tamaño. Las cuatro particiones de 2,5 MiB suman 10 MiB.

\textbf{Comprobación esperada.} La cota suma 10 MiB de entrada, 10 MiB de retorno de particiones, 10 MiB de entrega a W2 y 10 MiB de finales: 40 MiB. El observado también suma 40 MiB; agrupar no elimina las particiones distintas ni su relevo obligatorio.

\textbf{Resultado y decisión.} D02 y GD4 conservaron 40 MiB requeridos y observados, con DME igual a uno. Las trece entregas científicas permanecieron al pasar de cinco jobs a dos; como el observado ya coincide con la cota, agrupar consumidores de particiones distintas no redujo bytes. La ficha siguiente examina consumidores de un mismo archivo.

\subsection{Distribución de un archivo compartido}
\label{sec:prueba_b}
\textbf{Pregunta y propósito.} Cuando varios consumidores requieren un mismo archivo, \textquestiondown{}qué cambia al agruparlos? Se compara la reutilización con D02 y GD4, donde cada consumidor necesita una partición distinta.

\begin{table}[htbp]\centering\small
\caption{Configuraciones de la distribución de un archivo compartido}
\begin{tabularx}{\linewidth}{@{}P{1.3cm}YY@{}}
\toprule Caso & Configuración & Por qué se incluye\\\midrule
AC1 & Cuatro consumidores en jobs separados en W2; productor en W1. & Base con cuatro entradas del mismo archivo a W2.\\
AC2 & Dos jobs de consumo, uno para $C_1$ y $C_2$ y otro para $C_3$ y $C_4$. & Comprueba una entrada por grupo de dos consumidores.\\
AC4 & Un job de consumo para las cuatro tareas. & Comprueba si un grupo requiere una entrada.\\
\bottomrule
\end{tabularx}
\end{table}

\textbf{Factor y controles.} Varía el número de consumidores por job: uno, dos o cuatro. El productor permanece en W1 y los consumidores en W2. La entrada, $g$ y cada final miden 10 MiB; cada consumidor produce un final distinto que se entrega al coordinador.

\textbf{Comprobación esperada.} La cota es 70 MiB: 10 MiB de entrada a W1, 20 MiB por el retorno y la entrega de $g$, y 40 MiB de finales. El observado suma 60 MiB fijos (entrada, retorno de $g$ y finales) y una entrada de $g$ por job de consumo: 100, 80 o 70 MiB en total.

\textbf{Resultado y decisión.} AC1, AC2 y AC4 observaron 100, 80 y 70 MiB; DME fue 7/10, 7/8 y 1. AC4 alcanza la cota con una copia de $g$ a W2 y reduce el observado 30\,\% frente a AC1. La agrupación permite reutilizar $g$ dentro de cada job. En D02 y GD4, cada consumidor necesita una partición distinta, por lo que agrupar no elimina esas cuatro entradas.

\subsection{Redistribución}
\label{sec:prueba_r}
\textbf{Pregunta y propósito.} \textquestiondown{}Cómo se interpreta DME cuando una nueva asignación introduce destinos en varias etapas del mismo flujo? El flujo combina producción, agregación, división y consumo.

\begin{table}[htbp]\centering\small
\caption{Configuraciones del contraste de redistribución}
\begin{tabularx}{\linewidth}{@{}P{1.4cm}YY@{}}
\toprule Caso & Configuración & Por qué se incluye\\\midrule
R01-E & Cuatro productores, un agregador, un divisor y cuatro consumidores, todos en W1. Cada tarea tiene su job. & Es una base con tres clases de intermedios y todas las dependencias científicas locales en W1.\\
R02-E & Productores y divisor en W1. Agregador y consumidores en W2. Se conservan los diez jobs. & Introduce destinos entre máquinas en las tres etapas intermedias. Comprueba la referencia ante nuevas obligaciones.\\
\bottomrule
\end{tabularx}
\end{table}

\textbf{Factor y controles.} Se compara un factor de asignación de etapas con dos niveles completos. Cambian conjuntamente varias ubicaciones en R02-E. Se mantienen el grafo, las entradas y finales de 10 MiB, los cuatro primeros intermedios y las cuatro particiones de 10 MiB, el agregado de 40 MiB y la política de entrega. El cambio conjunto sirve para contrastar la suma por destino. Aislar efectos temporales requeriría otros contrastes.

\textbf{Comprobación esperada.} Las entradas y finales aportan 80 MiB. En R01-E no hay relevo intermedio obligatorio entre workers. En R02-E, los tres conjuntos intermedios aportan 40 MiB cada uno y requieren retorno y entrega: la cota es $80+2(40+40+40)=320$ MiB. El observado es 320 MiB en ambas configuraciones.

\textbf{Resultado y decisión.} R01-E tiene 80 MiB requeridos, 320 MiB observados y DME $=1/4$: existe exceso sobre la cota. R02-E tiene 320 MiB en ambos volúmenes y DME igual a uno: los recorridos intermedios son relevo obligatorio. No hay ahorro entre configuraciones; cambia la referencia por ubicación.

